\documentclass[10pt,a4paper]{amsart}
\usepackage[margin=19mm]{geometry}
\usepackage{amsmath,amssymb,mathtools}
\usepackage[scr=rsfs,cal=euler]{mathalfa}
\usepackage{xfrac}
\usepackage[unicode,hidelinks,bookmarks=false]{hyperref}
\newcommand{\ie}{i.e.\ }
\newcommand{\cf}{cf.\ }
\newcommand{\resp}{respectively\ }
\newcommand{\Subsection}{\subsection}
\providecommand{\nobreakdash}{\nobreak\hbox{-}}
\setlength{\emergencystretch}{2em}
\multlinegap0pt
\usepackage{bm}
\usepackage{bbm}
\usepackage{dsfont}
\usepackage{xspace}
\usepackage{cancel}
\usepackage{mathtools}
\usepackage{cleveref}
\usepackage{amsthm}
\makeatletter
\@ifundefined{lemma}{\newtheorem{lemma}{Lemma}}{}
\makeatother
\title[Uniqueness of bounded solutions to the fuzzy Landau and mult]{Uniqueness of bounded solutions to the fuzzy Landau and multiespecies Landau equations}
\author{F.-U. Caja-Lopez}
\date{Source: arXiv:2603.27509v1 (CC BY 4.0)}
\begin{document}
\maketitle

\noindent\emph{Source and licence.} Uniqueness of bounded solutions to the fuzzy Landau and multiespecies Landau equations. Version arXiv:2603.27509v1 (CC BY 4.0), distributed under the Creative Commons Attribution 4.0 International licence, as recorded in the arXiv licence metadata for this exact version and shown on the abstract page https://arxiv.org/abs/2603.27509v1. Full rights notice: licenses/arxiv-2603.27509-CC-BY-4.0.txt.\par\smallskip

\noindent\textit{Editorial scope.} Counted unit: the Lemma whose source label is \texttt{lem:commutators} in the article (source0.tex lines 2102--2116, source label \texttt{lem:commutators}), delivered together with the article's own complete original proof of it (source0.tex lines 2118--2233, one \texttt{proof} environment of 116 lines). No mathematics is written, repaired or re-derived: statement and proof are the article's own text line for line. Required context retained uncounted: lem:fuzzy\_commutator\_computation (lines 1812--1822); eq:analog\_heat\_multiespecies1 (lines 2057--2064); eq:analog\_heat\_multiespecies2 (lines 2057--2064). Every definition, display and label of those lines is retained unchanged. Omitted from this package and not redistributed: 1--1811, 1823--2056, 2065--2101, 2117--2117, 2234--2277. Nothing inside a retained excerpt is omitted or altered: every retained body is delivered exactly as the article's lines are written, so no omission receipt of any kind accompanies this package. Citations: the retained excerpts carry no citation command at all, so no bibliography entry of the article is transcribed and no reference is redistributed. Reference adaptation: none was needed and none was made; every reference inside the delivered unit resolves inside it, so no unresolved reference or citation is produced. Printed slips kept literal: no printed word, symbol, hypothesis or number of the counted statement or of its proof was altered. Layout and class: the article's own class is \texttt{amsart} with packages including amsmath, graphicx, amssymb, amsfonts, amsthm, todonotes, enumitem, array, float, dsfont, parskip, color, xcolor, geometry; the wrapper is the lane's pinned \texttt{amsart} editorial wrapper at 10pt on A4, which restates the theorem-like environments the retained text needs and adds the article's own macro definitions that the retained text needs (none), with extra wrapper packages (bm, bbm, dsfont, xspace, cancel, mathtools, cleveref). The delivered numbering is the unit's own. Assets: no retained span contains a figure environment, a graphics-inclusion command, a PDF-page inclusion, a table or a TikZ picture; no image file is copied into this package, the package declares no asset dependency and no image file is redistributed with it. Licence: version arXiv:2603.27509v1 of this article is distributed under the Creative Commons Attribution 4.0 International licence, as recorded in the arXiv licence metadata for this exact version and shown on the abstract page https://arxiv.org/abs/2603.27509v1. A full rights notice is delivered at licenses/arxiv-2603.27509-CC-BY-4.0.txt. Characters: every retained excerpt is pure ASCII, so the delivered engine, XeLaTeX with its default Latin Modern text font, reports no missing character.
\begin{lemma}\label{lem:fuzzy_commutator_computation}
  Let $b$ be one of the vector fields $\tilde b_j$ and $e\in \mathbb R^6\times \mathbb R^6$ a fixed vector. Then
  % 
  \begin{align*}
    I_{1}(e,b)+I_{2}(e,b) & 
    =-\int\left(e\cdot\nabla\left(b\cdot\nabla U\right)\right)^{2} r^\gamma \kappa P
    -\int\left([e,b]\cdot\nabla U\right)\left(e\cdot\nabla U\right)\left(b\cdot\nabla P\right) r^\gamma \kappa\\
    & +\int\left([e,b]\cdot\nabla\left(b\cdot\nabla U\right)\right)\left(e\cdot\nabla U\right) r^\gamma \kappa P
    +\int\left(e\cdot\nabla\left(b\cdot\nabla U\right)\right)\left([e,b]\cdot\nabla U\right) r^\gamma \kappa P.
  \end{align*}
\end{lemma}

\begin{align}
  2\frac{d}{dt}\sum_{i=1}^Nd_2^2(f^{(i)},g^{(i)}) & 
  =\sum_{i,j,k}c_{ij}\int\nabla\left(b_{k}^{i,j}\cdot\nabla\left(b_{k}^{i,j}\cdot\nabla U\right)\right)\cdot\nabla Ur_{ij}^{\gamma}P
  \label{eq:analog_heat_multiespecies1}\\
  & -\frac{1}{2}\sum_{i,j,k}c_{ij}\int\left(b_{k}^{i,j}\cdot\nabla\left(b_{k}^{i,j}\cdot\nabla|\nabla U|^{2}\right)\right)r_{ij}^{\gamma}P	
  \label{eq:analog_heat_multiespecies2}\\
  & +\frac{\gamma}{2}\sum_{i,j,k}c_{ij}\int\left(n^{i,j}\cdot\nabla U\right)\left(b_{k}^{i,j}\cdot\nabla\left(b_{k}^{i,j}\cdot\nabla U\right)\right)\,r_{ij}^{\gamma-2}P.\nonumber
\end{align}

\begin{lemma}\label{lem:commutators}
  Consider fixed $i,j,k$ then we have
  %
  \begin{equation*}
    (\ref{eq:analog_heat_multiespecies1})+(\ref{eq:analog_heat_multiespecies2}) \leq 
    -\frac{c_{ij}}{2} \int\left|\nabla\left(b_{k}^{i,j}\cdot\nabla U\right)\right|^{2}r_{ij}^{\gamma}P + C\max_{1\leq i\leq N}\left\Vert I_{\gamma}\rho_{i}\right\Vert _{L^{\infty}}\sum_{i=1}^Nd_2^2(f^{(i)},g^{(i)}),
  \end{equation*}
  %
  where $C$ only depends on $\min_i m_i$, $\max_{i,j}c_{i,j}$ and $I_\gamma$ denotes the following singular integral
  %
  \begin{equation}\label{eq:singular_integral}
    (I_{\gamma}\rho)(v):=\int_{\mathbb R^3}\left|v-v_*\right|^{\gamma}\rho(v_*)\,dv_*.
  \end{equation}
  %
\end{lemma}

\begin{proof}
  This is an adaptation of \cref{lem:fuzzy_commutator_computation} and the subsequent computations. This case is a little more involved since now the commutators don't all cancel out, so there are extra terms to be bounded. If we consider $i,j,k$ fixed and omit the coefficient $c_{ij}$ in front then we have
  %
  \begin{equation*}
    (\ref{eq:analog_heat_multiespecies1})  
    = \sum_{l,q}\left[I_{1}(b_{k}^{i,j},e_{l,q})+I_{1}(b_{k}^{i,j},\xi_{l,q})\right], \quad 
    (\ref{eq:analog_heat_multiespecies2}) 
    = \sum_{l,q}\left[I_{2}(b_{k}^{i,j},e_{l,q})+I_{2}(b_{k}^{i,j},\xi_{l,q})\right],
  \end{equation*}
  %
  where we defined
  %
  \begin{equation*}
    I_{1}(b,e)=\int\left(e\cdot\nabla\left(b\cdot\nabla\left(b\cdot\nabla U\right)\right)\right)\left(e\cdot\nabla U\right)r_{ij}^{\gamma}P,
    \quad I_{2}(b,e)=-\frac{1}{2}\int\left(b\cdot\nabla\left(b\cdot\nabla\left(e\cdot\nabla U\right)^{2}\right)\right)r_{ij}^{\gamma}P.
  \end{equation*}
  %
  Repeating the computations from Lemma \ref{lem:fuzzy_commutator_computation} we obtain again
  % 
  \begin{align*}
    I_{1}(b,e)+I_{2}(b,e) & =-\int\left(e\cdot\nabla\left(b\cdot\nabla U\right)\right)^{2}r_{ij}^{\gamma}P-\int\left([e,b]\cdot\nabla U\right)\left(e\cdot\nabla U\right)\left(b\cdot\nabla P\right)r_{ij}^{\gamma}	\\
    & +\int\left([e,b]\cdot\nabla\left(b\cdot\nabla U\right)\right)\left(e\cdot\nabla U\right)r_{ij}^{\gamma}P+\int\left(e\cdot\nabla\left(b\cdot\nabla U\right)\right)\left([e,b]\cdot\nabla U\right)r_{ij}^{\gamma}P.
  \end{align*}
  %
  %
  This yields the following expression:
  %
  \begin{align}
    (\ref{eq:analog_heat_multiespecies1})  + &(\ref{eq:analog_heat_multiespecies2})  =   
    -\int\left|\nabla\left(b_{k}^{i,j}\cdot\nabla U\right)\right|^{2}r_{ij}^{\gamma}P	\nonumber\\
    & -\sum_{l,q}\int\left[\left([e_{l,q},b_{k}^{i,j}]\cdot\nabla U\right)\left(e_{l,q}\cdot\nabla U\right)\left(b_{k}^{i,j}\cdot\nabla P\right)+\left([\xi_{l,q},b_{k}^{i,j}]\cdot\nabla U\right)\left(\xi_{l,q}\cdot\nabla U\right)\left(b_{k}^{i,j}\cdot\nabla P\right)\right]r_{ij}^{\gamma} 
    \raisetag{-6mm}\label{eq:commutators1}\\
    & +\sum_{l,q}\int\left[\left([e_{l,q},b_{k}^{i,j}]\cdot\nabla\left(b_{k}^{i,j}\cdot\nabla U\right)\right)\left(e_{l,q}\cdot\nabla U\right)+\left([\xi_{l,q},b_{k}^{i,j}]\cdot\nabla\left(b_{k}^{i,j}\cdot\nabla U\right)\right)\left(\xi_{l,q}\cdot\nabla U\right)\right]r_{ij}^{\gamma}P
    \raisetag{-6mm}\label{eq:commutators2}\\
    & +\sum_{l,q}\int\left[\left(e_{l,q}\cdot\nabla\left(b_{k}^{i,j}\cdot\nabla U\right)\right)\left([e_{l,q},b_{k}^{i,j}]\cdot\nabla U\right)+\left(\xi_{l,q}\cdot\nabla\left(b_{k}^{i,j}\cdot\nabla U\right)\right)\left([\xi_{l,q},b_{k}^{i,j}]\cdot\nabla U\right)\right]r_{ij}^{\gamma}P
    \raisetag{-6mm}\label{eq:commutators3}.
  \end{align} 
  %
  Let us now expand (\ref{eq:commutators1}) for $k=1$ as an example:
  %
  \begin{equation*}
    (\ref{eq:commutators1})=\alpha_{ij}\int\left[-\left(e_{i,2}\cdot\nabla U\right)\left(\xi_{j,3}\cdot\nabla U\right)\left(b_{1}^{i,j}\cdot\nabla P\right)+\left(e_{i,3}\cdot\nabla U\right)\left(\xi_{j,2}\cdot\nabla U\right)\left(b_{1}^{i,j}\cdot\nabla P\right)\right]r_{i,j}^{\gamma},
  \end{equation*}
  %
  % %
  % \begin{align*}
  %   (\ref{eq:commutators1}) & =-\int\Big(\overbrace{[e_{i,2},b_{1}^{i,j}]}^{\frac{e_{i,3}}{m_{i}}-\frac{\xi_{j,3}}{m_{j}}}\cdot\nabla u\Big)\left(e_{i,2}\cdot\nabla U\right)\left(b_{1}^{i,j}\cdot\nabla P\right)
  %   -\int\Big(\overbrace{[e_{i,3},b_{1}^{i,j}]}^{-\frac{e_{i,2}}{m_{i}}+\frac{\xi_{j,2}}{m_{j}}}\cdot\nabla U\Big)\left(e_{i,3}\cdot\nabla U\right)\left(b_{1}^{i,j}\cdot\nabla P\right)	\\
  %   & +\int\Big(\underbrace{[e_{i,2},b_{1}^{i,j}]}_{\frac{e_{i,3}}{m_{i}}-\frac{\xi_{j,3}}{m_{j}}}\cdot\nabla u\Big)\left(\xi_{j,2}\cdot\nabla U\right)\left(b_{1}^{i,j}\cdot\nabla P\right)+\int\Big(\underbrace{[e_{i,3},b_{1}^{i,j}]}_{-\frac{e_{i,2}}{m_{i}}+\frac{\xi_{j,2}}{m_{j}}}\cdot\nabla U\Big)\left(\xi_{j,3}\cdot\nabla U\right)\left(b_{1}^{i,j}\cdot\nabla P\right) \\
  %   & =\alpha_{ij}\int\left[-\left(e_{i,2}\cdot\nabla U\right)\left(\xi_{j,3}\cdot\nabla U\right)\left(b_{1}^{i,j}\cdot\nabla P\right)+\left(e_{i,3}\cdot\nabla U\right)\left(\xi_{j,2}\cdot\nabla U\right)\left(b_{1}^{i,j}\cdot\nabla P\right)\right]r_{i,j}^{\gamma}.
  % \end{align*}
  %
  where we have introduced the coefficients $\alpha_{ij} = \frac{1}{m_{i}}-\frac{1}{m_{j}}$. Integrating by parts and commuting then shows that for $k=1$
  %
  \begin{align*}
    (\ref{eq:commutators1}) & =  \alpha_{ij}\int\left(e_{i,2}\cdot\nabla\left(b_{1}^{i,j}\cdot\nabla U\right)\right)\left(\xi_{j,3}\cdot\nabla U\right)r_{ij}^{\gamma}P
    +\alpha_{ij}\int\left(\xi_{j,3}\cdot\nabla\left(b_{1}^{i,j}\cdot\nabla U\right)\right)\left(e_{i,2}\cdot\nabla U\right)r_{ij}^{\gamma}P	\\
    & -\alpha_{ij}\int\left(e_{i,3}\cdot\nabla\left(b_{1}^{i,j}\cdot\nabla U\right)\right)\left(\xi_{j,2}\cdot\nabla U\right)r_{ij}^{\gamma}P
    -\alpha_{ij}\int\left(\xi_{j,2}\cdot\nabla\left(b_{1}^{i,j}\cdot\nabla U\right)\right)\left(e_{i,3}\cdot\nabla U\right)r_{ij}^{\gamma}P\\
    & -\alpha_{ij}^{2}\int\left(e_{i,2}\cdot\nabla U\right)\left(\xi_{j,2}\cdot\nabla U\right)r_{ij}^{\gamma}P
    -\alpha_{ij}^{2}\int\left(e_{i,3}\cdot\nabla U\right)\left(\xi_{j,3}\cdot\nabla U\right)r_{ij}^{\gamma}P\\
    & -\alpha_{ij}^{2}\int\left(e_{i,2}\cdot\nabla U\right)^{2}r_{ij}^{\gamma}P
    -\alpha_{ij}^{2}\int\left(e_{i,3}\cdot\nabla U\right)r_{ij}^{\gamma}P.
  \end{align*}
  %
  All the previous terms are treated similarly using Young's inequality with a parameter $\varepsilon>0$. For example, for the first one we have
  %
  \begin{equation*}
    \alpha_{ij}\int\left(e_{i,2}\cdot\nabla\left(b_{1}^{i,j}\cdot\nabla U\right)\right)\left(\xi_{j,3}\cdot\nabla U\right)r_{ij}^{\gamma}P\leq\varepsilon\int\left|\nabla\left(b_{1}^{i,j}\cdot\nabla U\right)\right|^{2}r_{ij}^{\gamma}P+\frac{\alpha_{ij}^{2}}{4\varepsilon}\int\left(\xi_{j,3}\cdot\nabla U\right)^{2}r_{ij}^{\gamma}P.
  \end{equation*}
  %
  Now expand the integral remembering that $U=\bigoplus_{i}(u^{(i)}\oplus u^{(i)})$ and $P=\bigotimes_{i}(\rho^{(i)}\otimes\rho^{(i)})$
  %
  \begin{align*}
    \frac{\alpha_{ij}^{2}}{4\varepsilon}\int\left(\xi_{j,3}\cdot\nabla U\right)^{2}r_{ij}^{\gamma}P &
    =\frac{\alpha_{ij}^{2}}{4\varepsilon}\int\left(\partial_{3}u^{(j)}(v^{j}_*)\right)^{2}|v^{i}-v_*^{j}|^{\gamma}\rho^{(i)}(v^{i})\rho^{(j)}(v_*^{j})\,dv^{i}\,dv_*^{j}	\\
    % & =\frac{\alpha_{ij}^{2}}{4\varepsilon}\int\left(\partial_{3}u_{j}(w_{j})\right)^{2}\rho_{j}(w_{j})\left(I_{\gamma}\rho_{i}\right)(w_{j})\,dw_{j}\\
    & \leq\frac{\alpha_{ij}^{2}}{4\varepsilon}\Vert I_{\gamma}\rho^{(i)}\Vert _{L^{\infty}}d_{2}^{2}(f^{(j)},g^{(j)}),
  \end{align*}
  %
  which shows
  %
  \begin{align*}
    \alpha_{ij}\int\left(e_{i,2}\cdot\nabla\left(b_{1}^{i,j}\cdot\nabla U\right)\right)\left(\xi_{j,3}\cdot\nabla U\right)r_{ij}^{\gamma}P &
    \leq \varepsilon\int\left|\nabla\left(b_{1}^{i,j}\cdot\nabla U\right)\right|^{2}r_{ij}^{\gamma}P	\\
    & + \frac{\alpha_{ij}^{2}}{4\varepsilon}\Vert I_{\gamma}\rho^{(i)}\Vert _{L^{\infty}}d_{2}^{2}(f^{(j)},g^{(j)}).
  \end{align*}
  %
  The same estimate applies to the rest of the terms in the previous expansion of (\ref{eq:commutators1}) and for $k=2,3$. All in all, we obtain
  %
  \begin{equation*}
    (\ref{eq:commutators1}) \leq 6 \varepsilon\int\left|\nabla\left(b_{k}^{i,j}\cdot\nabla U\right)\right|^{2}r_{ij}^{\gamma}P
    +C\max_{1\leq i\leq N}\Vert I_{\gamma}\rho^{(i)}\Vert _{L^{\infty}}\sum_{i=1}^Nd_2^2(f^{(i)},g^{(i)}).
  \end{equation*}
  %
  The term $(\ref{eq:commutators2}) + (\ref{eq:commutators3})$ is bounded very similarly. 
  % , which we exemplify for $k=1$. After introducing the value of the commutators and some cancellation one obtains
  % %
  % \begin{align*}
  %   (\ref{eq:commutators2}) + (\ref{eq:commutators3}) &	
  %   =\alpha_{ij}\int\left(\xi_{j,3}\cdot\nabla\left(b_{1}^{i,j}\cdot\nabla U\right)\right)\left(e_{i,2}\cdot\nabla U\right)r_{ij}^{\gamma}P+\alpha_{ij}\int\left(e_{i,2}\cdot\nabla\left(b_{1}^{i,j}\cdot\nabla U\right)\right)\left(\xi_{j,3}\cdot\nabla U\right)r_{ij}^{\gamma}P\\
  %   & -\alpha_{ij}\int\left(\xi_{j,2}\cdot\nabla\left(b_{1}^{i,j}\cdot\nabla U\right)\right)\left(e_{i,2}\cdot\nabla U\right)r_{ij}^{\gamma}P-\alpha_{ij}\int\left(e_{i,3}\cdot\nabla\left(b_{1}^{i,j}\cdot\nabla U\right)\right)\left(\xi_{j,2}\cdot\nabla U\right)r_{ij}^{\gamma}P,
  % \end{align*}
  % %
  % which are bounded the same way as above. 
  All together, we have
  %
  \begin{equation*}
    (\ref{eq:analog_heat_multiespecies1})+(\ref{eq:analog_heat_multiespecies2}) 
    \leq
    -\left(1-10\varepsilon\right)\int\left|\nabla\left(b_{k}^{i,j}\cdot\nabla U\right)\right|^{2}r_{ij}^{\gamma}P
    +C_{\varepsilon}\max_{1\leq i\leq N}\Vert I_{\gamma}\rho^{(i)}\Vert _{L^{\infty}}\sum_{i=1}^Nd_2^2(f^{(i)},g^{(i)}).
  \end{equation*}
  %
  Picking $\varepsilon=\frac{1}{20}$ finishes the proof of the lemma after putting $c_{ij}$ in front again.
\end{proof}

\end{document}
